%% kapa: exam=1 round=19 year=2008 period=2 form=A subject=economics subject_ko=경제원론 seq=17 number=17
%% source: Q-Net 감정평가사 기출문제 (cst00309) artlSeq=1001962 file=extracted/1008/1001962/2008년 제19회 감정평가사 1차시험 기출문제/2교시.hwp
%% conversion: kordoc 4.24.51 -> tools/gen1.py (auto); review=proofread-sonnet; flags=none
\begin{question}{17}
변동환율제도 하에서 A국과 B국의 통화정책의 선택에 따라 나타나는 인플레이션율($\pi$)과 실업률($u$)의 변화분은 다음과 같다. 각국은 -△$\pi$/△$u$, 즉 실업률 1\% 상승에 따른 인플레이션율 감소분을 최대화하기를 원한다. 양국의 정책목표 및 정책수단 선택에 대한 설명으로 옳은 것은?

{\small\centering
\begin{tabular}{|c|c|c|c|}
\hline
\multicolumn{2}{|c|}{} & \multicolumn{2}{c|}{B국 통화정책} \\
\cline{3-4}
\multicolumn{2}{|c|}{} & 중립 & 긴축 \\
\hline
\multirow{4}{*}{\makecell{A국\\통화\\정책}} & \multirow{2}{*}{중립} & △$\pi^{A}$ = -1\%, △$u^{A}$ = 1\% & △$\pi^{A}$ = 0\%, △$u^{A}$ = 0.5\% \\
 & & △$\pi^{B}$ = -1\%, △$u^{B}$ = 1\% & △$\pi^{B}$ = -2\%, △$u^{B}$ = 1.75\% \\
\cline{2-4}
 & \multirow{2}{*}{긴축} & △$\pi^{A}$ = -2\%, △$u^{A}$ = 1.75 & △$\pi^{A}$ = -1.25\%, △$u^{A}$ = 1.5\% \\
 & & △$\pi^{B}$ = 0\%, △$u^{B}$ = 0.5\% & △$\pi^{B}$ = -1.25\%, △$u^{B}$ = 1.5\% \\
\hline
\end{tabular}\par}

\textsuperscript{*}윗첨자 $A,~B$는 각각 A국과 B국의 변수를 나타냄.
\begin{choices}
  \item A국은 인플레이션율을 그대로 유지하면서 실업률을 낮출 수 있다.
  \item B국은 인플레이션율을 낮추면서 실업률을 그대로 유지할 수 있다.
  \item B국이 긴축정책을 채택하면 A국은 중립정책을 채택하는 것이 좋다.
  \item 정책협조가 없다면 두 나라중 하나는 긴축정책을, 그리고 다른 하나는 중립정책을 채택하게 된다.
  \item A국과 B국이 모두 중립정책을 채택하면 파레토 효율적인 결과를 얻는다.
\end{choices}
\end{question}
\answer{5}
